\chapter{The Apparent Ontology}
\label{sec:false-ontology}
% ===================================================================

We can now state the main argument precisely, as a sketch pending
the filling of gaps identified in Section~\ref{sec:gaps}.

The cluster hierarchy --- the nerve of compact coherence clusters
and its recursive structure --- is not a product of the scientific
method.
It is an objective feature of the GSLT, arising from the
communication geometry of the population and the compactness
condition for consensus.
It is as real as the agents themselves.
The question is not whether the hierarchy exists but whether a
budget-limited agent running the scientific method will mistake
it for the whole story.
We argue that it will, and that this is not a failure of the
method but a consequence of the hierarchy being the dominant
experimentally accessible phenomenon at the agent's scale.

\section{What the Agent Sees}

Consider an agent $A$ embedded in a coherence cluster $\mathcal{C}_i$
at some level of the hierarchy.
$A$ runs the scientific method: it proposes HML formulae, tests them
with contexts, and updates its world model.
The method is sound: given unlimited budget, $A$ would converge
to the bisimulation quotient.
But $A$'s budget is finite, and three features of its situation
ensure that the cluster hierarchy is encountered first and
exhausts the budget before the deeper structure becomes accessible.

The experiments $A$ can run are bounded:
\begin{itemize}[itemsep=2pt]
  \item \emph{By budget}: $A$'s vectorial account $\vect{A}$ is finite.
        The bisimulation quotient is the limit of \emph{all} possible
        experiments; $A$ can run only finitely many.
  \item \emph{By coherence}: $A$ can only receive coherent responses
        from agents within $B_\epsilon(A)$.
        Agents beyond the coherence radius respond with noise, not
        with evidence.
  \item \emph{By the cluster boundary}: agents outside $\mathcal{C}_i$
        either do not respond coherently or, if they are in an
        overlapping cluster, respond only on the shared vocabulary
        of the overlap.
        The cluster boundary is not a wall the agent knows about;
        it simply marks where coherent evidence runs out.
\end{itemize}

\begin{proposition}[The Agent's World Model Converges to the Nerve]
\label{prop:false-ontology}
In the limit of $A$'s experimental budget, $A$'s world model
$\mathcal{W}$ converges to the fragment of the nerve $\nerve$
visible from $\mathcal{C}_i$: the sub-complex of $\nerve$
consisting of $\mathcal{C}_i$ and all clusters with which it shares
a non-empty intersection.
\end{proposition}

\begin{remark}
This is a sketch pending a precise definition of ``visible from
$\mathcal{C}_i$'' and a convergence proof.
The intuition is clear: $A$ can probe agents in $\mathcal{C}_i$
directly; it can probe agents in overlapping clusters through the
overlap vocabulary; agents in non-overlapping clusters are invisible.
The world model therefore reflects the nerve structure.
Crucially, the nerve structure is \emph{real} --- the clusters
genuinely exist and genuinely dominate the accessible evidence.
The agent is not hallucinating a structure that isn't there; it is
accurately theorizing a structure that is there, while being
unaware that the structure is itself composite.
\end{remark}

\section{Why the Deeper Ontology is Invisible}

The bisimulation quotient of $\GSLT$ --- together with the prime
combinator decomposition --- is not inaccessible because the
scientific method is defective.
It is inaccessible because it requires experimental resources that
exceed the budget available within any single coherence cluster.
Three compounding factors account for this:

\begin{enumerate}[label=\textbf{(\arabic*)}, itemsep=6pt]

  \item \textbf{Ontological isolation sets the terms of access.}
        $A$ can only observe encodings of other terms as they interact
        with it.
        The bisimulation quotient is the limit of \emph{all} possible
        experiments.
        With unlimited budget, the method converges there.
        With finite budget, it converges to whatever structure
        dominates the accessible experiments --- and that structure
        is the cluster hierarchy.

  \item \textbf{Communication degradation hides the global structure.}
        The agents whose behavior would reveal the global bisimulation
        quotient --- the ones that span cluster boundaries, connect
        distant parts of the population, or exhibit behavior only
        expressible in terms of the full GSLT --- are precisely the
        ones most distant in the communication geometry.
        Their responses arrive corrupted or not at all.
        The far reaches of the GSLT ontology are experimentally
        silent, not because they are absent but because the signal
        attenuates below the noise floor before it arrives.

  \item \textbf{Compactness makes the cluster self-confirming.}
        The coherence cluster $\mathcal{C}_i$ is compact, which means
        the agent's experiments converge to an internal consensus.
        Every hypothesis the agent proposes is tested against
        cluster-internal evidence, which confirms the cluster
        structure and does not probe beyond it.
        The cluster is not a prison; the agent could in principle
        ask questions whose answers require inter-cluster resources.
        But those experiments are expensive, the intra-cluster
        experiments are cheap, and a budget-limited agent will
        exhaust its resources on the cheap experiments first.
        By the time the cluster structure is confirmed and
        well-understood, the budget is gone.

\end{enumerate}

\section{Good Science, Limited Budget}

The agent $A$ is not making a mistake.
It is doing exactly what the scientific method prescribes: forming
the best theory consistent with the available evidence, within the
available resources.
The cluster hierarchy \emph{is} the available evidence.
It is a real structure, present in the GSLT independently of any
agent's theorizing.
The theory of the cluster hierarchy is correct at the scale it
describes.

What the agent lacks is not better method but more resources:
specifically, the ability to run experiments that cross the
coherence boundary, communicate with distant agents through the
noise, and probe the fine-grained bisimulation structure that
underlies the clusters.
With those resources --- with an unbounded budget --- the
scientific method would eventually push through to the bisimulation
quotient.
The agent is not epistemically defective; it is epistemically
bounded.

The hierarchy of cluster-level theories --- each level the
correct best theory at its scale, each level unaware of the
finer structure below --- is the computational analogue of the
hierarchy of effective field theories in physics.
Newtonian mechanics is correct at everyday scales.
Quantum field theory is correct at subatomic scales.
Neither is wrong; they are theories at different resolutions,
each the best available given the experimental resources of
their era.
The budget-limited agent's theory of the cluster hierarchy
stands in the same relation to the bisimulation quotient: correct
at its resolution, and silent about what lies below --- not
because what lies below is unreal, but because it is experimentally
unreachable within the budget.

% ===================================================================
